% Transcription — fonds Alexandre Grothendieck, shelfmark 53, batch 2 (pages 21-33).
% Established from the facsimile archives/batches/53/batch-02.pdf, cut from a
% copy of 53.pdf the user uploaded (PDF page k = archivists' page 20+k, no
% cover sheet), per the skill transcribe-grothendieck. The folder is
% thirty-three pages; this is the second and last of two batches. Batch 1
% (pages 1-20) was transcribed in parallel by another pass and had not been
% seen when this one was made, so the notation here is fixed from these
% pages alone and may need aligning with batch 1.
%
% Pass: Opus 5.5 (claude-opus-5-5), 2026-09-24 — first pass, unchecked against the pages by a human.
%
% Pages skipped: none. Pages summarised: none. All thirteen leaves (21-33)
% are manuscript notes in his hand, blue-black ink on ruled paper, fast
% drafting with many interlinear substitutions; none of them is the 1973
% letter named in the folder title, which must lie in batch 1 if anywhere.
% No date is written on these pages. He does not paginate this run; the
% leaves carry only the archivists' numbers. Several sentences continue
% from one leaf to the next (26 -> 27, 27 -> 28), which fixes their order.
%
% What the batch is about. Foundations of affine group schemes over a ring
% k, through coalgebras and distributions — the material of a course or a
% set of notes on algebraic groups. Page 21: a « Rectification » (an affine
% monoid of finite type over a field embeds as a closed multiplicative
% submonoid of M_n, not of GL_n), then Dist(X) = V(A) for an affine X,
% functoriality, Dist(X) x Dist(Y) -> Dist(X x Y), the point
% distributions X -> Dist(X), and X recovered from the coalgebra D as the
% x with Delta(x) = x (x) x. Page 22: Dist(G) for a monoid, actions of G
% versus actions of Dist(G), the hyperalgebra and the Lie algebra inside
% Dist(G, k), constant groups (group algebra) and diagonalisable groups
% (Dist(D(M))(k) = k^M with the product law). Pages 23-24: an action of an
% affine monoid G on a k-module M rewritten as u_0 : M -> M (x) A, the unit
% axiom (1') and the associativity axiom (2') as commutative diagrams, i.e.
% comodules. Pages 25-26: the functors V(M), W(M), full faithfulness,
% coalgebra structures on A as O-linear laws on V(A), comodules; the dual
% algebra A^ = V(A)(k) acting on M, and when A is projective of finite type
% (or a sum of such, with the topology of simple convergence) comodule
% structures equal continuous A^-module structures. Page 27: the
% Bil cont computation behind that, and stable quotients and submodules.
% Page 28: induced coaction on a submodule M'' -> M, uniqueness and
% existence when A is flat, M/M'' flat, or M'' a direct factor; a
% Proposition: when A is a sum of projectives of finite type, stable =
% sub-A^-module. Page 29: Corollaries 1-4 (stable hull, M the filtered
% union of its finite stable submodules, the same for bimodules, every
% coalgebra over a field is the filtered union of finite-dimensional ones).
% Page 30: Prop. — an affine group over a field is a filtered projective
% limit of affine groups of finite type, and one of finite type embeds in
% GL(n); proof by stable subalgebras and the antipode. Page 31: the algebra
% of distributions Dist(G)(k'), G(k') -> Dist(G)(k')^x, a cancelled example
% (finite constant group), and the antipode diagrams. Page 32: a list of
% examples (G_m, G_a, GL(n), SL(n), O(Q), Tr(n), Diag(n) ...) and a proof
% that W(M) : k' -> M (x) k' is representable iff M is projective of
% finite type. Page 33: V(M) represented by Sym*(M), the duality between
% V(M) and W(M), Aut_O(M) in End_O(M) representable for M projective of
% finite type, and a last line « fla, D(Gamma) ... ».
%
% Notation fixed for the batch: A the affine algebra, A^ (\check{A}) its
% dual V(A)(k), pi the comultiplication A -> A (x) A, epsilon the counit,
% sigma the antipode, u_0 (also omega_M, u_M, Delta_M) the coaction
% M -> M (x) A, O = O_k the structure ring functor. He writes the functor
% k' -> Hom_k(M, k') now V(M) now W(M) on page 25 (on page 32 W(M) is
% k' -> M (x) k'); the transcription follows the letter on the page and
% says so once. Dist is written Dist or Dis. His tensor product, ⊠ on
% page 21, is kept as \boxtimes.
%
% Hardest pages: 28 and 29, which are heavily reworked (whole lines struck,
% interlinear rewrites, a cancelled boxed remark); page 25's second margin
% note; page 26's cancelled paragraph. The formulas carry almost everywhere;
% connective words are often left \ill{}.
%
% The dating is the inventory's own for the folder (« 1973 », which is the
% letter's date); the manuscript notes are « s.d. » in the title, and these
% pages carry no date.

\documentclass[11pt,a4paper]{article}
% Le préambule commun à toutes les transcriptions du fonds Grothendieck.
%
% Il tient en une idée : **l'incertitude se compose, elle ne se résout pas.**
% Une transcription qui lisse un mot douteux a détruit la seule chose pour
% laquelle on la faisait — un lecteur doit pouvoir savoir, à chaque endroit,
% ce qui était sur le feuillet et ce qui a été ajouté. Les six macros qui
% suivent sont donc l'appareil critique, et non de la décoration.
%
% Les mêmes commandes sont reconnues par `scripts/render.mjs`, qui produit la
% vue de lecture du site. Une macro ajoutée ici doit l'être aussi là-bas,
% faute de quoi le rendu échouera bruyamment — ce qui est voulu : un rendu
% silencieusement faux est pire qu'un rendu absent.

\ProvidesPackage{grothendieck}[2026/08/08 Transcriptions du fonds Grothendieck]

\usepackage{amsmath,amssymb,amsthm}
\usepackage{mathrsfs}
\usepackage{stmaryrd}
% Les diagrammes commutatifs sont partout dans ce fonds : c'est la langue
% même de la géométrie algébrique de Grothendieck, et une transcription qui
% les rendrait en prose ne transcrirait rien.
\usepackage{tikz-cd}
% Français par défaut — c'est la langue des feuillets — mais l'anglais est
% chargé aussi : la traduction et le résumé sont des documents anglais, et la
% typographie française (espaces avant les deux-points) y serait une faute.
% Ces documents basculent par \selectlanguage{english} en tête de corps.
\usepackage[english,french]{babel}
\usepackage{fontspec}
\usepackage{geometry}
\geometry{a4paper,margin=25mm,marginparwidth=28mm}
\usepackage{marginnote}
\usepackage{xcolor}
% ulem, not soul. `soul`'s \st and \ul are built on a character-by-character
% scan that XeLaTeX's font machinery breaks: the document fails with an
% undefined control sequence rather than degrading. ulem does the same two jobs
% and is Unicode-engine safe. `normalem` stops it hijacking \emph, which the
% transcriptions use for Grothendieck's own emphasis.
\usepackage[normalem]{ulem}
\usepackage{hyperref}
\hypersetup{colorlinks=true,linkcolor=black,urlcolor=black,pdfborder={0 0 0}}

% --- Métadonnées du lot -----------------------------------------------------
% Reprises telles quelles par le script de rendu, qui les affiche en tête de
% la vue de lecture. Elles fixent ce que le fichier prétend couvrir : sans
% elles, une transcription flotte sans point d'attache dans un fonds de seize
% mille feuillets.

\newcommand{\folder}[1]{\gdef\grothFolder{#1}}
\newcommand{\batch}[1]{\gdef\grothBatch{#1}}
\newcommand{\leaves}[2]{\gdef\grothFirst{#1}\gdef\grothLast{#2}}
\newcommand{\foldertitle}[1]{\gdef\grothFoldertitle{#1}}
\gdef\grothFolder{?}\gdef\grothBatch{?}\gdef\grothFirst{?}\gdef\grothLast{?}\gdef\grothFoldertitle{}

% --- Le résumé --------------------------------------------------------------
%
% Il ouvre la lecture modernisée, sous le titre « Résumé », et il est composé
% en petit. Ce n'est pas un ornement : c'est le seul passage du document qui
% ne soit pas des mathématiques, et un lecteur doit pouvoir le reconnaître
% avant de l'avoir lu — puis le sauter s'il connaît déjà le sujet, ou s'y
% arrêter s'il ne le connaît pas. Le filet qui le suit marque la reprise du
% corps.

\newenvironment{resume}
  {\small\setlength{\parskip}{0.45em}}
  {\par\medskip\hrule\medskip}

% --- Mots-clés ---------------------------------------------------------------
%
% En anglais, en clôture du résumé de la lecture modernisée : le vocabulaire
% moderne sous lequel chercher ce que les feuillets construisent. Le manifeste
% du site (`npm run manifest`) les extrait pour étiqueter la cote entière —
% cette ligne est la source unique des tags, il n'y a pas de fichier à côté.
\newcommand{\keywords}[1]{\par\smallskip\noindent
  {\footnotesize\textsc{Keywords} --- \textit{#1}}\par}

% --- L'appareil critique ----------------------------------------------------

% Le numéro de feuillet, en marge. C'est l'ancre qui permet de régler un
% désaccord sur une formule en regardant *un* feuillet plutôt qu'une chemise
% de six cents. Le site s'en sert aussi pour tourner les pages du fac-similé
% quand on fait défiler la transcription.
\newcommand{\leaf}[1]{%
  \marginnote{\footnotesize\color{gray!70}\textsf{#1}}%
  \label{leaf:#1}\ignorespaces}

% Illisible : on ne devine pas. Un mot inventé qui se lit comme les autres est
% le pire résultat possible.
\newcommand{\ill}{\textcolor{red!60!black}{[\,\ldots\,]}}

% Lecture douteuse : le mot est donné, et signalé comme incertain.
\newcommand{\uncertain}[1]{\uline{#1}}

% Ajout de l'éditeur — une lettre manquante, un mot sous-entendu.
\newcommand{\add}[1]{\textcolor{blue!50!black}{[#1]}}

% Biffé par Grothendieck lui-même. Ce qu'il a rayé fait partie du document :
% une rature dit souvent où la pensée a changé de direction.
\newcommand{\struck}[1]{\sout{#1}}

% Note du transcripteur, sur l'état matériel du feuillet.
\newcommand{\note}[1]{\par\smallskip{\small\color{gray!60!black}\textit{[#1]}}\par\smallskip}

% Note marginale de Grothendieck, distincte du corps du texte.
\newcommand{\marginal}[1]{\par\smallskip\noindent\hspace{1em}%
  {\small\color{gray!60!black}#1}\par\smallskip}

% --- Titre ------------------------------------------------------------------

\AtBeginDocument{%
  \begin{center}
    {\large\bfseries Fonds Alexandre Grothendieck --- cote n\textsuperscript{o}~\grothFolder}\\[2pt]
    {\itshape\grothFoldertitle}\\[4pt]
    {\small feuillets \grothFirst--\grothLast\ (lot \grothBatch)}\\[2pt]
    {\footnotesize\color{gray!60!black}Transcription établie d'après le fac-similé numérisé
     par l'Université de Montpellier.\\
     Les lectures incertaines sont soulignées, les illisibles notées [\,\ldots\,],
     les ajouts entre crochets.}
  \end{center}
  \vspace{1em}}

\endinput


\folder{53}
\batch{2}
\pages{21}{33}
\dating{1973}
\watermark{Édition de démonstration}
\foldertitle{Groupes algébriques : notes manuscrites (s.d.), lettre (1973)}

\begin{document}

\page{21}
\emph{Rectification}\quad Si $G$ monoïde affine de t.f.\ sur un corps $k$, on a une immersion fermée multiplicative $G \to M_{n,k}$ (non $G \to \mathrm{GL}_{n,k}$), car on peut montrer que si $G$ est un monoïde de prés.\ finie sur un anneau $k$, tel que les translations à g (à dr.) \uncertain{sont} \ill{} \ill{} --- p.ex.\ s'il $\exists$ \uncertain{mon.} $G \to G'$, $G'$ un groupe --- alors $G$ est un objet \uncertain{gr.-pr.}).
\note{au-dessus de « groupe », un mot interlinéaire relié par un trait, peut-être « facteur… » ; la fin de la phrase, « objet gr.-pr. », est incertaine}

\uncertain{Soit} $X$ schéma affine sur $k$, \ill{} de l'algèbre $A$. On définit
\[
  \check{A} = \mathrm{Dist}(X)(k) ,
\]
\ill{} \ill{} le foncteur $k' \mapsto \mathrm{Dist}(X, k')$, noté
\[
  \mathrm{Dis}(X) \quad \bigl(\,= V(A)\ \text{comme Module sur } \mathcal{O}_k\bigr)
\]
\marginal{$\swarrow$ covariant}
\begin{enumerate}
  \item[a)] $\mathrm{Dist}(X)$ fonctoriel \uncertain{covariant} en $X$ ;
  \item[b)] $\mathrm{Dist}(X) \times \mathrm{Dist}(Y) \xrightarrow{\ \pi_{X,Y}\ } \mathrm{Dist}(X \times Y)$ bilinéaire
\end{enumerate}
[et universel pour les hom bilinéaires de $\mathrm{Dis}(X) \times \mathrm{Dis}(Y)$ dans un $V(M)$] i.e.\ \uncertain{définit} en fait un isom
\[
  \mathrm{Dist}(X) \boxtimes \mathrm{Dist}(Y) \xrightarrow{\ \sim\ } \mathrm{Dist}(X \times Y) .
\]
\note{la ligne a) porte, raturée au-dessus, « variant en $X$ » ou approchant ; la flèche marginale « covariant » y renvoie}
\begin{enumerate}
  \item[c)] On a un \uncertain{i}\ill{} fonctorielle
  \[
    X \xrightarrow{\ \alpha_X\ } \mathrm{Dist}(X) \qquad [\text{distribution ponctuelle}],
  \]
  compatible aux produits, i.e.\ …\ commutativité du
\end{enumerate}

\begin{tikzcd}
  X \times Y \arrow[r, "\alpha_X \times \alpha_Y"] \arrow[dr, "\pi_{X \times Y}"'] & \mathrm{Dist}(X) \times \mathrm{Dist}(Y) \arrow[d, "\pi_{X,Y}"] \\
  & \mathrm{Dist}(X \times Y)
\end{tikzcd}

\note{l'étiquette de la diagonale est écrite $\pi_{X\times Y}$ sur la page, alors que c'est la flèche ponctuelle de $X\times Y$ qu'appelle le diagramme ; la transcription garde la lettre de la page}

\struck{Si $G$ est muni d'une loi de comp., $\mathrm{Dist}(G)$ aussi. Elle est unitaire ssi celle de $\mathrm{Dist}(G)$ l'est.}
\note{deux lignes biffées d'un trait, puis de traits verticaux}
En fait on peut reconstruire $X$ à partir de $\mathrm{Dist}(X) = D$ et la structure \uncertain{co-algébrique}
\[
  D \xrightarrow{\ \Delta\ } D \boxtimes D \qquad \bigl[V(A) \boxtimes V(B) \overset{\mathrm{def}}{\simeq} V(A \otimes B)\bigr]
\]
\[
  [\text{NB } \boxtimes = \otimes \text{ si } A \text{ ou } B \text{ proj.\ de t.f.}]
\]
doit être une « application diagonale » ass.\ comm.\ et counitaire [i.e.\ muni de $D \xrightarrow{\ \varepsilon\ } \mathcal{O} \ldots$] comme formé des $x \in D(k')$ tels que
\[
  \boxed{\ \Delta(x) = x \boxtimes x\ }
\]
\marginal{NB on \uncertain{réinterprète} $D$ \uncertain{comme} $\mathrm{Dist}$ en \ill{} \ill{} $X$ : muni de \uncertain{diagonale} commutative, counitaire}
\note{note marginale en oblique, dans le coin inférieur gauche, lue en partie ; « $\mathrm{Dist}(X)$ » y paraît sous une forme abrégée}

\page{22}
[cet exprime que l'hom $A \to k'$ est un hom d'algèbres]

Si $G$ est muni d'une loi de comp., \ill{} \ill{} $\mathrm{Dist}(G)$, celle-ci est (ass.\ comm.) ssi celle de $G$ l'est, et \struck{unitaire} $G$ admet une unité bilatère ssi $\mathrm{Dist}(G)$ a une unité \struck{\ill{}} $e$ telle que $\delta(e) = e \boxtimes e$.

Dire que $G$ opère sur un module $M$ \add{(g.)} signifie que $\mathrm{Dist}(G)$ y opère à gauche --- et cette opération étend celle de $G$.
\note{« (g.) » est inscrit au-dessus de la ligne, après $M$ ; l'étendue de la reprise « et cette opération étend » est de lecture probable, non sûre}

\textbf{NB} Dans certains cas, $\mathrm{Dist}(X)$ se reconstitue à partir de
\[
  \mathrm{Dist}(X, k) = \mathrm{Dist}(X)(k),
\]
p.ex.\ si $A$ est \struck{\ill{}} projectif de t.f.\ sur $k$, ou plus gén.\ si $A$ est somme de modules projectifs de t.f.\ (p.ex.\ $k$ un corps) --- mais alors il faut considérer $\mathrm{Dist}(X, k)$ comme module topologique.

\textbf{Re NB} \uncertain{Surtout} \uncertain{remarquons} pour mémoire. En plus de $G$, $\mathrm{Dist}(G)$ [voire $\mathrm{Dist}(G, k)$] contient toute l'information infinitésimale de $G$ : p.ex.\ \emph{l'algèbre de Lie} de $G$ est contenue dans $\mathrm{Dist}(G, k)$, \struck{et la \ill{}} sa loi de $[\ ]$ est déduite par le crochet de l'algèbre, et on \uncertain{a} aussi sa puiss.\ $p$-ième aussi --- de sorte que l'hyperalgèbre.
\note{la phrase s'arrête sur « l'hyperalgèbre. » ; « Surtout remarquons » est une lecture très incertaine de la première ligne}

Cas des schémas constants : on retrouve \emph{l'algèbre} d'un groupe fini.

Cas des gr.\ \uncertain{pr.} diagonalisables
\[
  A = k(M) \quad \text{algèbre du groupe à coeff.\ dans } k
\]
\[
  \begin{array}{lcl}
    A \otimes A \longrightarrow A & \text{provient de} & M \times M \longrightarrow M \ \text{loi de groupe} \\
    A \otimes A \longleftarrow A & \text{---} & M \times M \longleftarrow M \ \text{diagonale}
  \end{array}
\]
\[
  \mathrm{Dist}(D(M))(k) = k^M
\]
avec la loi d'algèbre \emph{produit} (ne dépend pas de la str.\ \uncertain{de groupe} de $M$ !) [Combinaison \uncertain{linéaire} infinie de caractères].
\note{sous « $A\otimes A \leftarrow A$ », un signe « $=$ » au crayon, peu marqué ; « diagonale » est écrit au-dessus de la seconde ligne}

\page{23}
$A$ algèbre d'un $k$-\struck{groupe} monoïde affine $G$

$M$ $k$-module

opération de $G$ sur $M$ : \struck{\ill{}}
\[
  \begin{cases} 1)\ 1x = x \\ 2)\ g(g'x) = gg'x \end{cases}
  \qquad \text{i.e.}\quad G \longrightarrow \underline{\mathrm{End}}_k(M) = \underline{\mathrm{End}}_{\mathcal{O}}(W(M))
\]
« Opération sans axiomes » correspond à \struck{\ill{}}
\[
  u : M_A \xrightarrow[A\text{-lin}]{} M_A, \qquad M_A = A \otimes_k M,
\]
ou
\[
  u_0 : M \longrightarrow A \otimes_k M \qquad k\text{-lin.}
\]
\note{deux traits courbes, sur $A\otimes_k M$, semblent indiquer l'échange des facteurs, en $M\otimes_k A$ ; c'est l'ordre qu'emploient les formules qui suivent}

\struck{L'axiome $1x = x$}

Si $g \in G(k')$ correspond à un hom de $k$-algèbres $A \xrightarrow{g} k'$, alors $g_M$ sera l'endomorphisme \struck{\ill{}} de $(k', g) \otimes_A M$ \struck{défini} déduit \struck{de $g \otimes_A i$} par changement de base, i.e. l'unique hom $M_{k'} \to M_{k'}$ tel que le diagramme d'hom $A$-linéaires

\begin{tikzcd}
  M_A \arrow[r, "u"] \arrow[d, "g \otimes \mathrm{id}_M"'] & M_A \arrow[d, "g \otimes \mathrm{id}_M"] \\
  M_{k'} \arrow[r, "g_M"'] & M_{k'}
\end{tikzcd}

\note{à gauche du carré, entre crochets, $M$ envoie deux flèches, vers $M_A$ et vers $M_{k'}$}
est commutatif ; mais cela signifie aussi que le diagramme

\begin{tikzcd}
  M \arrow[r, "u_0"] \arrow[d, "\mathrm{can}"'] & M \otimes A \arrow[d, "\mathrm{id}_M \otimes g"] \\
  M \otimes k' \arrow[r, "g_M"'] & M \otimes k'
\end{tikzcd}

est commutatif, ou encore, notant que la donnée de $g_M$ est équivalente à celle de $g_M^{\circ} = g_M \circ \mathrm{can} : M \to M \otimes k'$, à la \struck{f\ill{}} commutativité de

\begin{tikzcd}
  M \arrow[r, "u_0"] \arrow[dr, "g^{\circ}_M"'] & M \otimes A \arrow[d, "\mathrm{id}_M \otimes g"] \\
  & M \otimes k'
\end{tikzcd}

Pour obtenir l'axiome 1), on fait $k' = k$, $g = \struck{\ill{}}(\varepsilon : A \to k)$, d'où

\textbf{(1')}

\begin{tikzcd}
  M \arrow[r, "u_0"] \arrow[dr, "="'] & M \otimes A \arrow[d, "\mathrm{id}_M \otimes \varepsilon"] \\
  & M
\end{tikzcd}

comm.\ i.e.\ $(\mathrm{id}_M \otimes \varepsilon)\, u_0 = \mathrm{id}_M$.

\page{24}
Pour écrire 2), on \struck{écrit} prend $k' = A \otimes A$, et on choisit les deux hom
\[
  A \mathrel{\substack{\xrightarrow{\ g\ } \\ \xrightarrow[\ g'\ ]{}}} A \otimes A
\]
\struck{\ill{}} donnés par
\[
  g(\lambda) = \lambda \otimes 1, \qquad g'(\lambda) = 1 \otimes \lambda ,
\]
et on écrit l'égalité comme une égalité de deux endom.\ de \struck{$M_{A\otimes A}$} $M_{k'} = M_{A\otimes A}$ qu'on va écrire --- puis on écrira l'égalité \add{de leurs composés} avec $M \xrightarrow{\mathrm{can}} M_{A\otimes A}$ --- donc finalement il faudra définir deux $k$-hom $M \to M_{A \otimes A}$, savoir
\[
  x \longmapsto g\,x', \qquad x \longmapsto g'\,x'
\]
(où $x'$ est l'image de $x$ dans $M_{k'} = M_{A\otimes A}$).
\note{l'« égalité de leurs composés » est une lecture de sens : la ligne porte « l'égalité » suivi d'un mot bref illisible}

Or
\[
  (gg')_M\, x' = (gg')^{\circ}_M\, x = (\mathrm{id}_M \otimes gg') \circ u_0 :
\]

\begin{tikzcd}
  M \arrow[r, "u_0"] \arrow[dr, "(gg')^{\circ}_M"'] & M \otimes A \arrow[d, "\mathrm{id}_M \otimes \pi"] \\
  & M \otimes A \otimes A
\end{tikzcd}

\note{à droite de ce diagramme, une règle de calcul sur les éléments, du type $y \otimes \lambda \mapsto \ldots$, est raturée et illisible ; $gg'$ n'est autre que la comultiplication $\pi : A \to A \otimes A$}
\[
  g_M(g'_M x') = \ ?
\]
Calculons $g'_M x' \in M_{k'}$ ($= (g'_M)^{\circ} x$), c'est $(\mathrm{id}_M \otimes g')(u_0(x))$ :

\begin{tikzcd}
  M \arrow[r, "u_0"] & M \otimes A \arrow[d] \\
  & M \otimes A \otimes A
\end{tikzcd}

$y \otimes \lambda \longmapsto y \otimes 1 \otimes \lambda$.

Donc $g_M(g'_M x')$ s'obtient en appliquant au résultat $g_M$, donné par
\[
  g_M(y \otimes \lambda \otimes \mu) = g^{\circ}_M(y)\, \lambda \otimes \mu = (\mathrm{id}_M \otimes g)(u_0(y))\, \lambda \otimes \mu ,
\]
\uncertain{en} \ill{} donc
\[
  g_M(y \otimes 1 \otimes \lambda) = \underbrace{(\mathrm{id}_M \otimes g)\, u_0(y)}_{\in\, M \otimes A \otimes 1}\, 1 \otimes \lambda ,
\]
donc

\begin{tikzcd}
  M \arrow[d, "u_0"'] & \\
  M \otimes A \arrow[r, "u_0 \otimes \mathrm{id}_A"] & M \otimes A \otimes A
\end{tikzcd}

\note{c'est le diagramme de $(g_M g'_M)^{\circ}$ ; dans $M \otimes A \otimes A$ au bout de la flèche, il souligne d'un trait ondulé le facteur $M \otimes A$ ; une ligne raturée court sous le cadre}
Donc l'axiome 2 s'exprime comme

\textbf{(2')}

\begin{tikzcd}
  M \arrow[r, "u_0"] \arrow[d, "u_0"'] & M \otimes A \arrow[d, "\mathrm{id}_M \otimes \pi"] \\
  M \otimes A \arrow[r, "u_0 \otimes \mathrm{id}_A"'] & M \otimes A \otimes A
\end{tikzcd}

\emph{commutatif}.

\page{25}
\struck{\textbf{NB} L'homom.\ $A \to A \otimes A$ s'interprète comme un hom \ill{}}
\note{ligne biffée d'un trait, et en outre barrée de grandes boucles ; au-dessus, « de $k$-modules », relié par un trait, semble pris dans la même rature}

\textbf{NB} À tout $k$-module $M$ associons le foncteur en $\mathcal{O}$-modules
\[
  W(M) : k' \longmapsto \mathrm{Hom}_{k\text{-}\mathrm{Mod}}(M, k')
\]
qui dépend de façon contravariante de $M$. On \uncertain{trouve} :
\begin{enumerate}
  \item[a)] $M \mapsto V(M)$, $\quad \mathrm{Mod}(k)^{\circ} \longrightarrow \mathrm{Mod}(\mathcal{O})$ est pl.\ fidèle, i.e.
  \[
    \forall M, N \qquad \mathrm{Hom}_k(M, N) \xrightarrow[\sim]{\ \alpha_{M,N}\ } \mathrm{Hom}_{\mathcal{O}}(W(N), W(M))
  \]
  \struck{\ill{}}
  \note{sous la formule, une ligne entière est raturée et illisible ; l'étiquette de l'alinéa, « b) », y est d'abord biffée puis reprise à la ligne suivante}
  \item[b)] $\forall M, N$ \struck{\ill{}} l'accouplement canonique
  \[
    V(M) \otimes_{\mathcal{O}} V(N) \longrightarrow V(M \otimes_k N) ,
  \]
  \ill{} pas un isom, car il n'est pas vrai que
  \[
    \bigl[\mathrm{Hom}_k(M, k') \otimes_{k'} \mathrm{Hom}_{k'}(N', k') \simeq \mathrm{Hom}_{k'}(M' \otimes_{k'} N', k')\bigr]
  \]
  \note{interligne, au-dessus : $M'^{\vee} \otimes_{k'} N'^{\vee} \to (M' \otimes_k N')^{\vee}$}
  définit cependant, en prenant les Hom dans \ill{} un $V(P)$, un hom fonctoriel
  \[
    \mathrm{Hom}_{\mathcal{O}}(V(M \otimes_k N), V(P)) \xrightarrow{\ \sim\ } \mathrm{Hom}_{\mathcal{O}}(V(M) \otimes_{\mathcal{O}} V(N), V(P))
  \]
  avec $\mathrm{Hom}_{\mathcal{O}}(V(M\otimes_k N), V(P)) \simeq \mathrm{Hom}_k(P, M \otimes N)$.
\end{enumerate}
\note{il écrit tantôt $V$, tantôt $W$ pour le même foncteur ; la page ne permet pas de les distinguer, et la transcription suit la lettre qu'il trace}

On trouve donc ainsi
\begin{itemize}
  \item[$\cdot$)] Une loi de composition $\mathcal{O}$-linéaire sur un $V(A)$ ($A$ un $k$-module) \uncertain{est} une comultiplication $A \xrightarrow{\pi} A \otimes A$.
  Une \uncertain{unité} \ill{} pour la loi de composition est une (counité) $A \xrightarrow{\varepsilon} k$.
  \note{« augmentation » est écrit au-dessus de « counité », que cerne un trait}
  Une \ill{} la loi de composition est \ill{} (comm.) ssi la loi d'algèbre est \ill{} (comm.).
  \note{sens attendu : la loi de composition est associative (commutative) ssi la loi de coalgèbre est coassociative (cocommutative) ; les mots eux-mêmes ne se lisent pas}
  Si $M$ est un module, un accouplement $\mathcal{O}$-bilinéaire
  \[
    V(A) \times V(M) \longrightarrow V(M)
  \]
  est une \ill{} \uncertain{diagonale} $M \to A \otimes M$. Si $A$ muni de $\varepsilon$ \uncertain{et} $\pi$ est unitaire, alors \struck{la loi de} l'accouplement est unitaire ssi le diagramme (1') \uncertain{comm.} ; \struck{Si $A$ muni d'une} \ill{} \quad $\lambda(\mu x) = (\lambda\mu)x$ \quad (ass.) ssi le diagramme 2' est commutatif.
\end{itemize}
\marginal{on dit aussi : $M$ comodule}
\marginal{\ill{} dit aussi : \ill{} comodule, \ill{} \ill{} \ill{} signifie que $V(A)$ opère sur $W(M)$ à gauche}
\note{deux notes cerclées dans la marge gauche, reliées par un trait à « module » et à l'accouplement ; la seconde, écrite en oblique et serrée, ne se lit qu'en partie}

\page{26}
Soit donc $A$ une coalgèbre (ass.\ unit.) et $M$ un comodule \add{à g.} sur $A$, la coalgèbre $\check{A} = V(A)(k)$ opère à droite sur $V(M)$, donc à gauche sur $W(M)$, ou ce qui revient au même, sur $M$, donc \struck{$V(M$} $M$) devient un $\check{A}$-\struck{alg} module. Si
\[
  \xi : A \longrightarrow k \quad (k\text{-lin}) \in \check{A}, \qquad x \in M,
\]
on
\[
  \xi \cdot x = (\mathrm{id}_M \otimes \xi)(u_0(x)),
\]
i.e.
\[
  \xi_M = (\mathrm{id}_M \otimes \xi) \circ u_0 .
\]
\note{« à g. » est inscrit au-dessus de « comodule » ; la lecture « à gauche » est suggérée par la phrase même}
Notons que si $A$ est \struck{\ill{}} projectif de t.f.\ sur $k$, \struck{on a un isom.,} on a un isom \struck{$M \otimes A \simeq \mathrm{Hom\,cont}_k(\ldots)$}
\[
  M \otimes A \simeq \mathrm{Hom}_k(\check{A}, M) ,
\]
donc la donnée des $\bigl[M \xrightarrow{u_0} M \otimes_k A\bigr]$ [satisfaisant certaines conditions] \struck{\ill{}} équivaut à la donnée de
\[
  \check{A} \longrightarrow \mathrm{Hom}_k(M, M)
\]
satisfaisant des conditions pareilles. On trouve ainsi que la donnée d'une structure de $A$-comodule \struck{\ill{}} sur $M$ \ill{}, que si le $k$-module $A$ est projectif de t.f.\ $A$, la donnée d'une structure de coalgèbre (ass.\ unit.) équivaut à celle d'une str.\ d'algèbre (ass.\ unit.) sur $\check{A}$.

\struck{Si $k$ est un corps, on a \ill{} \ill{} $A = \varinjlim A_i$, $A_i$ les \ill{} \ill{} de dim.\ finie, d'où $\check{A} = \varprojlim \check{A}_i$ (lim.\ projective) \ill{} de $\check{A}$} \marginal{gén.\ si $A$ est somme directe de modules de t.f.\ proj.} Munissant $\check{A}$ de la topologie de la convergence simple, on a
\[
  M \otimes_k A \simeq \mathrm{Hom\,cont}_k(\check{A}, M),
\]
\note{le passage « Si $k$ est un corps, … » est biffé par plusieurs traits obliques et horizontaux ; seules des bribes s'en lisent. La note marginale semble remplacer l'hypothèse biffée}
donc
\[
  \mathrm{Hom}(N, M \otimes_k A) \simeq \mathrm{Hom\,cont}_k\bigl(\check{A}, \underbrace{\mathrm{Hom}_k(N, M)}_{\text{muni de la top.\ de la conv.\ simple}}\bigr)
\]
donc on trouve qu'une structure de comodule (ass.\ unit.) sur $M$ revient à la donnée d'une
\note{la phrase se poursuit à la page suivante}

\page{27}
structure de $\check{A}$-module sur $M$ tel que pour tout $x \in M$, $\xi \mapsto \xi x$ soit continue (i.e.\ s'annule sur un des $J_i \subset \check{A}$ …). De même, on voit que la donnée d'une structure de coalgèbre (co.\ unit.) sur le $k$-module $A$ \ill{} équivaut à la donnée d'une \struck{opération} structure de $k$-algèbre (ass.\ unitaire) sur $\check{A}$ qui soit \emph{continue}
[car si $\check{P}, \check{Q}, \check{R}$ sont \ill{} \ill{}, \struck{\ill{} continues}
\note{au-dessus de $\check{P}$, $\check{Q}$, $\check{R}$ : $k^{\Sigma}$, $k^{J}$, $k^{\Lambda}$, reliés par $\simeq$ ou $=$ ; le deuxième exposant est incertain. Au-dessus du mot biffé, « bil. »}
\struck{$\check{P} \times \check{Q} \longrightarrow \check{R} = \varprojlim \check{R}_\gamma$}

\begin{align*}
  \mathrm{Bil\,cont}_k(\check{P}, \check{Q}; \check{R})
  &\simeq \mathrm{Bil\,cont}_k\bigl(\check{P}, \check{Q}; \varprojlim_{\gamma} \check{R}_\gamma\bigr) \\
  &\simeq \varprojlim_{\gamma} \bigl[\mathrm{Bil\,cont}_k(\check{P}, \check{Q}; \check{R}_\gamma)\bigr] \\
  &\simeq \varprojlim_{\gamma} \Bigl[\varinjlim_{\alpha,\beta} \underbrace{\mathrm{Bil}_k(\check{P}_\alpha, \check{Q}_\beta; \check{R}_\gamma)}_{\mathrm{Hom}_k(\check{P}_\alpha \otimes \check{Q}_\beta,\, \check{R}_\gamma)}\Bigr] \\
  &\simeq \varprojlim_{\gamma} \varinjlim_{\alpha,\beta} \mathrm{Hom}_k(R_\gamma, P_\alpha \otimes Q_\beta) \\
  &\simeq \varprojlim_{\gamma} \mathrm{Hom}_k(R_\gamma, P \otimes Q) \\
  &\simeq \mathrm{Hom}_k(R, P \otimes Q) \qquad !\ ]
\end{align*}
\note{une ligne « $\simeq$ » laissée vide et biffée entre la troisième et la quatrième égalité ; dans la première ligne, « cont » est surchargé sous « Bil »}

On est dans le cas envisagé \struck{\ill{}} si $A$ est limite inductive filtrante \struck{On a également les \ill{} \ill{}} \emph{\uncertain{dénombrable}} de sous-modules projectifs \add{de t.f.}, à morphismes de transition \struck{ayant des} \struck{supplém.} \uncertain{mon.}\ \uncertain{univ.}\ [i.e.\ ayant des \uncertain{rétractions} à g.]

Mais dans le cas général, on ne peut oublier $V(A)$ \add{et ses opérations} au profit de $\check{A} = V(A)(k)$ …
\note{« et ses opérations » est écrit au-dessus de la ligne et appelé par un trait}

\emph{Modules quotients et sous-modules stables par une co-opération de la coalgèbre $A$} : Soit $M \twoheadrightarrow N'$ un ép., alors
\[
  V(N') \subset V(M) \quad \text{mono},
\]
et il est clair ce que signifie que $V(N')$ est stable ; on dira alors que $N'$ est stable par la \uncertain{coop.}\ de $M$. Cela signifie aussi qu'$\exists$ un hom $N' \xrightarrow{\ \omega_{N'}\ } N' \otimes A$ rendant commutatif

\begin{tikzcd}
  M \arrow[r, "\omega_M"] \arrow[d] & M \otimes A \arrow[d, "g \otimes \mathrm{id}_A"] \\
  N' \arrow[r, "\omega_{N'}"'] & N' \otimes A
\end{tikzcd}

et cet $\omega_{N'}$ est unique et satisfait aux axiomes d'unitarité et associativité. Si $A$ est l'alg.\ affine
\note{l'épimorphisme est tracé $M \to N'$, surmonté d'une lettre peu nette, qui doit être le $g$ de la flèche verticale ; au bas de la page, « $N'$ » est inscrit sous « cet » ; la phrase se poursuit à la page suivante}

\page{28}
d'un monoïde affine $G$, cela signifie aussi qu'il existe une opération de $G$ sur $M'$ \struck{que} telle que $M \to M'$ soit comp.\ \ill{} op., et il est clair (comme $W(M') \to W(M)$ \uncertain{inj.}) que cette opération \uncertain{sera} unique.
\note{un mot interlinéaire, peut-être « devient », est écrit sous $M'$ et appelé par un trait vers la parenthèse}

Si \ill{} un mono
\[
  M'' \overset{i}{\hookrightarrow} M
\]
alors il ne sera pas vrai en général que $V(M) \to V(M'')$ soit surjectif, ni que $W(M'') \to W(M)$ soit injectif, donc \ill{} cette condition \ill{} \ill{} \ill{} à parler d'opération induite, \struck{\ill{}} On dira qu'il \emph{existe} une op.\ induite si on peut faire opérer $G$ sur $M''$ de telle façon que $M'' \to M$ soit comp.\ aux opérations. Mais il n'est plus clair que l'op.\ « induite » soit unique ! En termes d'alg.\ comm., la question est de trouver un $\Delta_{M''} : M'' \to M'' \otimes A$ rendant commutatif

\begin{tikzcd}
  M'' \arrow[r, "u_{M''}\,?"] \arrow[d, hook, "i"'] & M'' \otimes A \arrow[d, "i \otimes \mathrm{id}_A"] \\
  M \arrow[r, "u_M"'] & M \otimes A
\end{tikzcd}

On peut dire que c'est unique \add{\uncertain{et existe}} si $i \otimes \mathrm{id}_A = i_A$ soit un mono, et si [$i \otimes \mathrm{id}_{A \otimes A}$ l'est aussi, \ill{} \ill{} \ill{} \ill{} ; alors $M'$ est stable ssi $M''$ l'est].
\note{passage interlinéaire serré et raturé en partie, à droite : on y lit encore « $\mathrm{Im}(M'') \subset \mathrm{Im}\,M''_{\ldots}$ », « associativité » et, souligné, « existe » ; le reste est illisible}

\struck{Ceci} Ces conditions \uncertain{envisagées} sont ok dans les cas suivants :
\begin{enumerate}
  \item[a)] $A$ est plat (v.\ \ill{} le cas \uncertain{courant})
  \item[b)] $M/M''$ est $A$-plat
  \item[c)] $M''$ facteur direct dans $M$.
\end{enumerate}
\struck{\ill{}}
\note{à droite de a), trois lignes biffées de grands zigzags, qui finissaient sur « dans $M''\otimes A$ » ; en dessous, cerclé : « Alors l'intersection \uncertain{finie} de sous-modules stables est stable, kif kif sommes quelconques »}
\marginal{\uncertain{supposons} ces conditions, \ill{} $M'$ stable ssi $M''$ stable}
\note{la note marginale, en oblique et biffée de traits obliques, est reliée par une accolade à la liste a)--c)}

Dans le cas $A$ libre [ou plus gén.\ somme de modules proj.\ de t.f.] on voit que $M'$ est stable ssi $M'$ est stable \struck{\ill{} sous l'action de $\check{A}$}, i.e.\ \ill{} un sous-$\check{A}$-module. On trouve donc

\emph{Proposition}\quad Si $A$ est somme de modules proj.\ de t.f., alors un sous-$k$-module d'un $A$-comodule \uncertain{$M$} est stable ssi il \struck{est stable sous \ill{}} est un sous-$\check{A}$-module.

\page{29}
\emph{Corollaire 1}\quad Pour tout \struck{\ill{}} sous-module $M'$ de $M$, il y a un plus petit sous-module \add{$M'^{S}$} stable de $M$ qui contienne $M'$, et c'est aussi le sous $\check{A}$-module engendré. Si $M'$ est de t.f.\ \struck{sur} \add{\uncertain{sur}} $k$, $M'^{S}$ aussi, donc \ill{}
\note{« $M'^{S}$ » est inscrit au-dessus de « sous-module » ; l'exposant $S$ est une lecture de la lettre tracée, qui ressemble aussi à un $G$}

\emph{Corollaire 2}\quad $M$ est limite inductive filtrante de ses sous-modules de t.f.\ \emph{stables}.
\note{la ligne se termine par un mot bref biffé}

\struck{Corollaire 3} \struck{\ill{}} Considérons $A$ comme bimodule sur lui-même \add{\uncertain{\ill{}}} à g.\ et à droite, \struck{alors} donc comme bimodule sur $\check{A}$, \struck{\ill{} \ill{}, lorsque} \struck{alors \ill{} \ill{} $H$ \ill{}}, considérons les sous-bimodules $B \subset A$. \struck{On \ill{} que} \marginal{Analogue Cor 1, 2} \struck{\ill{} \ill{} \ill{}} \struck{$A$ \ill{} $A$ considéré comme \ill{}} $A$ \add{la coalgèbre produit ($A^{\circ}$ l'opposée de $A$)} \struck{$A$} \struck{sous} \ill{} $A \otimes A^{\circ}$ et on module \ill{} celle-ci, on trouve directement, on trouve
\note{ce passage est remanié à l'extrême : un titre « Corollaire 3 » biffé, plusieurs amorces raturées, un ajout interlinéaire « la coalgèbre produit ($A^{\circ}$ l'opposée de $A$) » relié par un long trait ; la syntaxe exacte n'est pas restituable}

\emph{Cor 3}\quad \struck{\ill{}} Tout sous-bimodule $B$ de $A$ est contenu dans un plus petit sous-bimodule $B^{S}$ \struck{stable par} stable à g.\ et à droite (i.e.\ le sous-$\check{A}$-bimodule engendré), et si $B$ est de t.f.\ sur $k$, $B^{S}$ aussi. Donc $A = \varinjlim A_i$, les $A_i$ étant stables à g.\ et à droite.

Or cette stabilité \add{pour un $B$ donné} signifie que $\pi : A \to A \otimes A$ envoie $B$ dans $(B \otimes A) \cap (A \otimes B)$. Si \struck{$B$ \ill{}} $A/B$ est plat sur $k$ (\uncertain{par ex.} le cas \uncertain{courant}) cela signifie aussi que $\pi$ envoie $B$ dans $B \otimes B$, donc veut dire qu'il y a une structure de coalgèbre induite. On trouve donc

\emph{Cor.\ 4}\quad Toute coalgèbre \add{$A$} sur un corps est $\varinjlim$ \add{fil.} de coalgèbres de t.f., donc $\check{A}$ est $\varprojlim$ \add{fil.\ stricte} d'algèbres finies sur $k$.
\note{« fil. » et « fil. stricte » sont des ajouts interlinéaires ; la première limite est inductive, la seconde, soulignée d'une flèche vers la gauche, projective}

\page{30}
\emph{Prop.}\quad $G$ \struck{\ill{}} \add{groupe (monoïde)} affine sur le corps $k$. Alors
\begin{enumerate}
  \item[a)] $G$ \struck{\ill{}} est limite projective \add{\uncertain{filtrante}, \uncertain{stricte}} de \add{groupes} \struck{\ill{}} (\uncertain{monoïdes}) affines de t.f.\ sur $k$, \struck{\ill{} des morphismes \ill{}}
  \item[b)] Si $G$ est un monoïde \add{(groupe)} affine de t.f.\ sur $k$, on peut trouver un \uncertain{entier} $n$ et une immersion fermée
  \[
    G \longrightarrow \mathrm{M}(n)_k \qquad (\text{resp.\ } G \to \mathrm{GL}(n)_k)
  \]
  \uncertain{soumise} aux multiplications.
\end{enumerate}
\note{« entier » est lu d'après le sens, le mot étant tracé d'un seul long trait. Le $\mathrm{M}$ est surchargé sur un $\mathrm{GL}$ ; en exposant, une petite croix}

\emph{Dém.}\quad $G \longleftrightarrow A$, \struck{\ill{}} et $A = \varinjlim_i A_i$, les $A_i$ les sous-algèbres de t.f.\ de $A$. On s'intéresse aux sous-algèbres $B$ qui sont « stables » par $\pi : A \to A \otimes A$, et il est clair que si $M \subset A$ est stable sous $\pi$, la \uncertain{sous-algèbre} \add{$B$} (qu'elle engendre) aussi. Donc le corollaire 2 nous dit que l'on peut supposer les $A_i$ stables par $\pi$. Si $G$ un groupe, d'où antipodisme $\sigma$ de $A$, alors si $M$ stable \uncertain{sous} $\sigma$, alors l'algèbre $B$ aussi. Or si $M$ stable par $\pi$, $\sigma M$ aussi [$M$ sous-module à g.\ \add{(à dr.)} \uncertain{sous} $\check{A}$ $\Longleftrightarrow$ $\sigma M$ sous-module à dr.\ (g.)\ \uncertain{sous} $\check{A}$], donc $M + \sigma M$ aussi, donc on peut supposer les $A_i$ stables par $\sigma$, donc les $G_i$ sont des groupes.

Supposons maintenant $A$ de t.f., donc on a un sous $k$-module \add{de t.f.} $M$ de $A$ qui engendre $A$ comme \struck{\ill{}} algèbre. \struck{Considérons} On peut supposer $M$ stable par $\check{A}$ à \struck{\uncertain{gauche}} droite, donc que $G$ y opère à \struck{droite} \ill{}. Donc
\[
  G \longrightarrow \struck{\ill{}}\ \underline{\mathrm{Aut}}_k(M) \quad (\simeq \mathrm{GL}(n)_k)
\]
où $n = \dim M$. Prouvons que c'est un mono. On a \struck{\ill{} \ill{}} \struck{$g \mapsto \varepsilon(f)$} …\ Considérons \add{pour $\varphi \in M$} la fonction sur $G$
\[
  g \longmapsto \varepsilon(g \cdot \varphi) = (g\varphi)(1) = \varphi(g)
\]
\note{au-dessus de « la fonction », « $\varphi \in M$ » ; sous $g \cdot \varphi$, un mot interlinéaire, peut-être « élém. »}
\struck{\ill{}} Donc $M \subset \mathrm{Im}\bigl[\text{anneau \uncertain{affine} de } \underline{\mathrm{End}}_k(M) \to A\bigr]$
donc $V$ \ill{} affine de $\underline{\mathrm{End}}_k(M)$ et \ill{} fonction \uncertain{coordonnée} de $\mathrm{GL}(n)$ \uncertain{s'envoient} \emph{sur} $A$, ok.

\page{31}
Si $A$ est l'algèbre \add{affine} d'un monoïde \add{affine $G$} sur $k$ [i.e.\ $A$ est muni, en plus de la comultiplication $\pi$, d'une multiplication \add{ass.\ et unitaire} bilinéaire \struck{qui en fait une} \add{ass.\ unitaire} qui est de plus \emph{commutative}, et $\pi$ est un hom d'algèbres] alors $\forall k'$, les éléments de \add{$V(A)(k') =$} $\mathrm{Hom}(A, k')$ \add{i.e.\ les hom.\ \uncertain{linéaires} $A \to k'$} \struck{\ill{}} s'appellent les « distributions de $G$ à \struck{\ill{}} coeff \add{\uncertain{h-mod.}} dans $k'$ », notés $\mathrm{Dist}(G)(k')$ [\add{algèbre des distributions}]. Donc \struck{on a} elles forment une $k'$-algèbre [« l'algèbre du \struck{groupe} \add{monoïde} $G$ à coeff.\ dans $k'$ »] et on a un homomorphisme de monoïdes
\[
  G(k') \hookrightarrow \mathrm{Dist}(G)(k')^{\times}
\]
et si $G$ est un groupe
\[
  G(k') \hookrightarrow \mathrm{Dist}(G)(k')^{*} \qquad [\text{pratiquement jamais un isom}]
\]
\note{« ass.\ et unitaire » est écrit au-dessus de « multiplication », « ass.\ unitaire » au-dessus des mots biffés ; « h-mod. », au-dessus de « coeff », est très incertain}

\struck{Ex : $G$ un groupe fini \struck{constant} ordinaire $\Gamma$, $G = \Gamma_k$, \ill{} \ill{}, $A =$ algèbre des fonctions \add{à valeurs dans $k$} sur $\Gamma$ \add{\uncertain{postulat}}}
\note{l'exemple est encadré puis biffé de quatre traits obliques}

\begin{tikzcd}
  G \arrow[r, "\mathrm{id}_G \times \sigma"] \arrow[dr, "p"'] & G \times G \arrow[r, "\pi"] & G \\
  & e \arrow[ur, "\varepsilon"'] &
\end{tikzcd}

\note{au-dessus, « $g\,g^{-1} = 1$ » ; la pointe de la première flèche est peu nette, le sens retenu est celui qu'impose le diagramme dual}
\struck{$A \to A \otimes A$}

\begin{tikzcd}
  A & A \otimes A \arrow[l] & A \arrow[l, "\pi"'] \arrow[dl, "\varepsilon"] \\
  & k \arrow[ul, "p"] &
\end{tikzcd}

\note{sur la flèche de gauche : $x\,\sigma y \longleftarrow\!\shortmid\ (x, y)$}
\marginal{$\pi(x) = $ \quad $\mathrm{Dist}(X) \xrightarrow{\ \pi\ } \mathrm{Dist}(X) \times \mathrm{Dist}(X)$, \quad $\varepsilon$ \ill{} $\mathrm{Dist}(X)$ \ill{} $e$}
\marginal{$X \times Y \ldots A \otimes A,\ x \otimes y \ldots$ \ill{}}
\note{deux notes marginales écrites le long du bord gauche, tournées d'un quart de tour ; la seconde est une esquisse de flèches entre $X\times Y$, $A\otimes A$ et $k$, qu'on ne restitue pas}

\page{32}
\emph{Exemples de groupes algébriques affines sur $k$.}
\[
  \begin{array}{llll}
    \mathbb{G}_m,\ \mathbb{G}_a, & \mathrm{GL}(n), & O(Q), & \text{s-groupe de } \mathrm{GL}(n) \text{ invariant des tenseurs} \\
    & \mathrm{SL}(n) & \ \cap & \quad \ldots \\
    & & \mathrm{GL}(n) & \\
    & \mathrm{Tr}(n) & & \\
    & \mathrm{Tr}_0(n) & & \\
    \mathrm{Diag}_0(n) \subset \mathrm{Diag}(n) & & & \\
    \quad \wr & & & \\
    \mathbb{G}_m & & &
  \end{array}
\]
\note{à gauche de $\mathrm{Tr}_0(n)$, une abréviation brève, peut-être « Sp. » ; l'isomorphisme vertical relie $\mathrm{Diag}_0(n)$ à $\mathbb{G}_m$}
immersion de $\mathbb{G}_a$ dans $\mathrm{GL}(2)$

Soit $M$ un $k$-module.
\[
  W(M) : k' \longmapsto M \otimes_k k' \qquad (\text{un faisceau } \mathcal{O}_k\text{-Module, pas seulement un groupe})
\]
Pour que ce soit représentable, il faut que $W(M)$ transforme injections $k' \to k''$ en injections, donc que $M$ soit \emph{plat} sur $k$. \struck{\ill{}} Il faut qu'il transforme produits (infinis) d'algèbres en produits
\[
  M \otimes_k \prod k'_i \longrightarrow \prod_i M \otimes_k k'_i \quad \text{bijectif.}
\]
\struck{On sait que ceci} faisant tous les $k'_i$ égaux à $k$, on trouve
\[
  M \otimes k^I \xrightarrow{\ \sim\ } M^I .
\]
Et on sait tout de suite que cela implique que $M$ est de type fini, donc \struck{\ill{}} quotient d'un libre \struck{fini} \struck{$\to R \to L$}
\[
  0 \to R \to L \to M \to 0
\]
\struck{\ill{}} et comme $M$ est plat, cela reste exacte par $\otimes_k k^I$, donc $R \otimes k^I \xrightarrow{\sim} R^I$, donc $R$ de type fini, donc $M$ de présentation finie, donc $M$ plat de prés.\ finie, donc $M$ projectif de t.f.\ (facteur direct d'un libre).

Inversement, si $M$ projectif de t.f.\ i.e.\ facteur direct d'un libre $L = k^n$, par un projecteur $p$, on a
\[
  M \otimes_k k' = \mathrm{Ker}\bigl((1 - p)_{k'} : k'^n \to k'^n\bigr)
\]
\[
  W(M) = \mathrm{Ker}\bigl(E^n \xrightarrow{\ 1-p\ } E^n\bigr)
\]
c'est repr.
\note{$E$ désigne sans doute le schéma affine standard $\mathbb{G}_a$ ; la page ne le définit pas}

\page{33}
\[
  V(M) = \mathrm{Hom}_{k\text{-mod}}(M, k') = \mathrm{Hom}_{k'\text{-mod}}(M \otimes_k k', k')
\]
\struck{est} représentable par $\mathrm{Sym}^{*}_k(M) = \sum_{n \geqslant 0} \mathrm{Sym}^{n}_k(M)$
\[
  \Bigl[\ V(M) \times W(M) \longrightarrow \mathcal{O}_k = W(k) \qquad V(M) \xrightarrow{\ \sim\ } \underline{\mathrm{Hom}}_{\mathcal{O}_k}\bigl(W(M), W(k)\bigr)
\]
\note{le second membre est expliqué par une accolade : $k' \longmapsto \mathrm{Hom}\bigl(W(M)_{k'}, W(k')\bigr)$, avec $W(M)_{k'} \simeq W(M_{k'})$ au-dessus, et en dessous $= \mathrm{Hom}_{k'}(M_{k'}, k') = \mathrm{Hom}_k(M, k')$. Le crochet ouvert n'est pas refermé}
\struck{Or} \quad $M \longmapsto W(M)$
\[
  \mathrm{Mod}(k) \longrightarrow \mathrm{Mod}(\mathcal{O}_k) \quad \text{est \emph{pleinement fidèle}.}
\]
\struck{Mod} \quad $\underline{M} \longmapsto V(M)$ \qquad idem après ext.\ de la base à $k'$
\struck{De \ill{}} \struck{\ill{}} Donc on a une \uncertain{dualité} parfaite entre $V(M)$ et $W(M)$ …
\note{« Donc on a une … » est encadré à gauche d'un trait, sous deux mots biffés ; « dualité » est proposé pour un mot bref commençant par « b » ou « d », peut-être « bi-dualité »}
\[
  \underline{\mathrm{Aut}}_{\mathcal{O}}(M) \subset \underline{\mathrm{End}}_{\mathcal{O}}(M) \qquad \text{représentable si } M \text{ est proj.\ t.f.}
\]

\uncertain{fla}, $D(\Gamma)$ …
\note{dernière ligne, sous un trait horizontal : « fla » (ou « fl$_a$ ») suivi de $D(\Gamma)$ et de points de suspension ; la lecture du premier mot est incertaine}

\end{document}
